% ============================================================================
%  Mono-age disk populations and the age dependence of the J_R-h_Z relation
%  ApJ -- AASTeX v7
%  Skeleton created 2026-10-26.
% ============================================================================
\documentclass[twocolumn]{aastex701}

\usepackage{amsmath}
\usepackage{amssymb}

\newcommand{\JR}{J_R}
\newcommand{\hZ}{h_Z}
\newcommand{\sigz}{\sigma_z}
\newcommand{\sigR}{\sigma_R}
\newcommand{\Lsun}{L_\odot}
\newcommand{\kpc}{\,\mathrm{kpc}}

\shorttitle{Mono-age populations and the J_R-h_Z relation}
\shortauthors{Vera-Ciro}

\begin{document}

% Working title, to be settled once all seven Auriga halos are analyzed.
\title{The Age Dependence of the Spiral $\JR$--$\hZ$ Relation}

\author{Carlos Vera-Ciro}
\affiliation{Independent Researcher}
\email{caverac@gmail.com}

\begin{abstract}
% TODO: write last.
\end{abstract}

\keywords{Galaxy dynamics (591) --- Galaxy structure (622) --- Milky Way dynamics
(1051) --- Stellar dynamics (1596)}

\section{Introduction}
\label{sec:intro}
% TODO.

\section{Simulations and data}
\label{sec:data}
% TODO. The citation policy of the Auriga data release asks every publication that uses
% its data to cite the overview paper (Grand2017) and the data release paper (Grand2024).

\section{Method}
\label{sec:method}
% TODO.

\section{Mono-age populations in Auriga}
\label{sec:auriga}
% TODO.

\section{The Milky Way}
\label{sec:milkyway}
% TODO.

\section{The $\JR$--$\hZ$ relation}
\label{sec:relation}
% TODO.

\section{Discussion and conclusions}
\label{sec:conclusions}
% TODO.

\begin{acknowledgments}
We have used simulations from the Auriga Project public data release \citep{Grand2024}
available at \url{https://wwwmpa.mpa-garching.mpg.de/auriga/data.html}.

This work has made use of data from the European Space Agency (ESA) mission
{\it Gaia} (\url{https://www.cosmos.esa.int/gaia}), processed by the {\it Gaia}
Data Processing and Analysis Consortium (DPAC,
\url{https://www.cosmos.esa.int/web/gaia/dpac/consortium}). Funding for the DPAC
has been provided by national institutions, in particular the institutions
participating in the {\it Gaia} Multilateral Agreement
\citep{Gaia2016,Gaia2023}.

Funding for the Sloan Digital Sky Survey IV has been provided by the Alfred P. Sloan
Foundation, the U.S. Department of Energy Office of Science, and the Participating
Institutions. SDSS acknowledges support and resources from the Center for
High-Performance Computing at the University of Utah. The SDSS web site is
www.sdss4.org. SDSS is managed by the Astrophysical Research Consortium for the
Participating Institutions of the SDSS Collaboration including the Brazilian
Participation Group, the Carnegie Institution for Science, Carnegie Mellon University,
Center for Astrophysics $|$ Harvard \& Smithsonian (CfA), the Chilean Participation
Group, the French Participation Group, Instituto de Astrof\'{i}sica de Canarias, The
Johns Hopkins University, Kavli Institute for the Physics and Mathematics of the
Universe (IPMU) / University of Tokyo, the Korean Participation Group, Lawrence Berkeley
National Laboratory, Leibniz Institut f\"{u}r Astrophysik Potsdam (AIP),
Max-Planck-Institut f\"{u}r Astronomie (MPIA Heidelberg), Max-Planck-Institut f\"{u}r
Astrophysik (MPA Garching), Max-Planck-Institut f\"{u}r Extraterrestrische Physik (MPE),
National Astronomical Observatories of China, New Mexico State University, New York
University, University of Notre Dame, Observat\'{o}rio Nacional / MCTI, The Ohio State
University, Pennsylvania State University, Shanghai Astronomical Observatory, United
Kingdom Participation Group, Universidad Nacional Aut\'{o}noma de M\'{e}xico,
University of Arizona, University of Colorado Boulder, University of Oxford, University
of Portsmouth, University of Utah, University of Virginia, University of Washington,
University of Wisconsin, Vanderbilt University, and Yale University. We use APOGEE DR17
\citep{Blanton2017,Majewski2017,Abdurrouf2022,Wilson2019,Gunn2006}, with the astroNN
\citep{LeungBovy2019a,LeungBovy2019b,Mackereth2019} and DistMass
\citep{StoneMartinez2024} value-added catalogs.

This publication makes use of data products from the Two Micron All Sky Survey, which
is a joint project of the University of Massachusetts and the Infrared Processing and
Analysis Center/California Institute of Technology, funded by the National Aeronautics
and Space Administration and the National Science Foundation \citep{Skrutskie2006}.

Guoshoujing Telescope (the Large Sky Area Multi-Object Fiber Spectroscopic Telescope
LAMOST) is a National Major Scientific Project built by the Chinese Academy of Sciences.
Funding for the project has been provided by the National Development and Reform
Commission. LAMOST is operated and managed by the National Astronomical Observatories,
Chinese Academy of Sciences. We use the LAMOST-Gaia red-clump catalog of
\citet{Huang2020}.

We use the density fits of \citet{Imig2025}, released under CC BY 4.0 at
\url{https://github.com/astrojimig/mw_density_imig2025}.
\end{acknowledgments}

% Remove once the text cites its references.
\nocite{*}

\bibliographystyle{aasjournalv7}
\bibliography{ms}

\end{document}
